\documentclass[10pt]{article}
\usepackage[T1]{fontenc}
\usepackage[utf8]{inputenc}
\usepackage[french]{babel}
\usepackage{amsmath,amssymb,mathrsfs,mathpazo}
\usepackage{geometry}
\geometry{a4paper,margin=2cm}
\usepackage{graphicx,array,tabularx,tabularray,colortbl,multirow,multicol,booktabs,enumitem,pifont,ulem,fancybox,diagbox,cancel}
\usepackage{pstricks,pstricks-add,pst-plot,pst-tree,pst-node,pst-text,pst-func,pst-math,pst-bspline,pst-3dplot,pst-eucl,pst-solides3d}
\let\pscalculate\undefined
\usepackage{pst-bezier}
\usepackage{tikz,pgfplots,tkz-tab}
\usetikzlibrary{arrows,arrows.meta,patterns,calc,positioning,trees,decorations.markings,shapes.geometric,shapes.misc,shapes.arrows,intersections,decorations.pathmorphing,decorations.pathreplacing,backgrounds,shadows,babel}
\pgfplotsset{compat=1.18}
\usepgfplotslibrary{fillbetween}
\usepackage[np]{numprint}
\usepackage{listings,makecell,pas-tableur,twoopt}
\let\barre\undefined
\usepackage[pstricks]{tablvar}
\definecolor{darkgreen}{RGB}{0,100,0}
\providecommand{\red}{\color{red}}
\providecommand{\blue}{\color{blue}}
\providecommand{\green}{\color{green}}
\providecommand{\black}{\color{black}}
\providecommand{\white}{\color{white}}
\providecommand{\gray}{\color{gray}}
\providecommand{\cyan}{\color{cyan}}
\providecommand{\magenta}{\color{magenta}}
\providecommand{\yellow}{\color{yellow}}
\providecommand{\violet}{\color{violet}}
\providecommand{\orange}{\color{orange}}
\providecommand{\euro}{\texteuro}
\providecommand{\EUR}{\texteuro}
\providecommand{\ieme}{\textsuperscript{e}}
\providecommand{\degres}{\ensuremath{^\circ}}
\providecommand{\degree}{\ensuremath{^\circ}}
\providecommand{\vv}[1]{\overrightarrow{#1}}
\providecommand{\Coord}[1]{#1}
\providecommand{\decosix}{}
\providecommand{\starredbullet}{\ensuremath{\bullet}}
\pagestyle{empty}
\setlength{\parindent}{0pt}
\usepackage{adjustbox}\DeclareUnicodeCharacter{2236}{:}\DeclareUnicodeCharacter{2190}{\ensuremath{\leftarrow}}
\providecommand{\R}{\mathbb{R}}
\providecommand{\N}{\mathbb{N}}
\providecommand{\Z}{\mathbb{Z}}
\providecommand{\Q}{\mathbb{Q}}
\providecommand{\D}{\mathbb{D}}
\providecommand{\C}{\mathbb{C}}
\providecommand{\e}{{\mathrm{e}}}
\providecommand{\vect}[1]{\overrightarrow{#1}}
\providecommand{\vectt}[1]{\overrightarrow{\mathrm{#1}}}
\providecommand{\vvt}[1]{\overrightarrow{\mathrm{#1}}}
\providecommand{\np}[2][]{{#2~#1}}
\providecommand{\vv}[1]{{\overrightarrow{#1}}}
\renewcommand{\d}{{\mathrm{d}}}
\providecommand{\degres}{{^\circ}}
\providecommand{\degre}{{^\circ}}
\providecommand{\blue}{\color{blue}}
\providecommand{\red}{\color{red}}
\providecommand{\enskip}{\quad}
\providecommand{\llbracket}{\lbrack\!\lbrack}
\providecommand{\rrbracket}{\rbrack\!\rbrack}
\providecommand{\euro}{{\text{€}}}
\providecommand{\no}{{\text{n°}}}
\providecommand{\interval}[2]{{[#1;#2]}}
\providecommand{\intervFO}[2]{{[#1;#2[}}
\providecommand{\dotfill}{\ldots}

\providecommand{\ds}{\displaystyle}
\providecommand{\textding}[1]{#1}
\providecommand{\Oij}{\(\mathrm{O};\vec{\imath},\vec{\jmath}\)}
\providecommand{\Oijk}{\(\mathrm{O};\vec{\imath},\vec{\jmath},\vec{k}\)}
\providecommand{\Ouv}{\(\mathrm{O};\vec{u},\vec{v}\)}
\providecommand{\intervFF}[2]{\ensuremath{[#1;#2]}}
\renewcommand{\euro}{\ifmmode\text{\texteuro}\else\texteuro\fi}
\DeclareUnicodeCharacter{00B0}{\ifmmode^{\circ}\else\textdegree\fi}
\DeclareUnicodeCharacter{20AC}{\ifmmode\text{\texteuro}\else\texteuro\fi}
% Énoncés : Éducation nationale. Transcription et corrigés : APMEP, Denis Vergès et contributeurs. Découpage et classement : CoopMaths, Sébastien Lozano, Éric Elter et contributeurs. Adaptation de la présentation : Maths Explorer.
% CC BY-SA : https://coopmaths.fr/www/about/
% https://coopmaths.fr/alea/static/bac/2024/tex/bac_2024_05_sujet2_ameriquenord_4.tex
\begin{document}
\noindent\textit{APMEP / CoopMaths. CC BY-SA. Présentation adaptée par Maths Explorer.}\par\medskip
Soit $a$ un réel strictement positif.

On considère la fonction $f$ définie sur l'intervalle $]0~;~+\infty[$ par 

\[f(x) = a \ln (x).\]

On note $\mathcal{C}$ sa courbe représentative dans un repère orthonormé.

Soit $x_{0}$ un réel strictement supérieur à 1.

\medskip

\begin{enumerate}
\item Déterminer l'abscisse du point d'intersection de la courbe $\mathcal{C}$ et de l'axe des abscisses.
\item Vérifier que la fonction $F$ définie par $F(x) =a [x \ln (x) - x]$ est une primitive de la fonction $f$ sur l'intervalle $]0~;~+\infty[$.
\item En déduire l'aire du domaine bleuté en fonction de $a$ et de $x_{0}$.

\end{enumerate}

\begin{center}
\psset{unit=2cm,arrowsize=2pt 3}
\begin{pspicture*}(-1,-1)(4.5,2)
\psaxes[linewidth=1.25pt,Dx=5,Dy=5]{->}(0,0)(-1,-1)(4.5,2)
\pscustom[fillstyle=solid,fillcolor=blue]{\psplot[plotpoints=2000,linewidth=1.25pt,linecolor=blue]{1}{1.8}{x ln 1.3 mul}\psline(1.8,0.764)(1.8,0)(1,0)}
\psplot[plotpoints=2000,linewidth=1.25pt,linecolor=blue]{0.1}{4.5}{x ln 1.3 mul}
\uput[dr](1,0){1}\uput[d](1.8,0){$x_0$}\uput[ul](1.8,0.764){$M$}
\uput[dr](4,1.6){\blue $\mathcal{C}$}\uput[dl](0,0){O}
\uput[d](4.4,0){$x$}\uput[l](0,1.9){$y$}
\end{pspicture*}
\end{center}

On note $T$ la tangente à la courbe $\mathcal{C}$ au point $M$ d'abscisse $x_{0}$.

On appelle $A$ le point d'intersection de la tangente $T$ avec l'axe des ordonnées et $B$ le projeté orthogonal de $M$ sur l'axe des ordonnées.

\begin{center}
\psset{unit=2cm,arrowsize=2pt 3}
\begin{pspicture*}(-1,-1)(4.5,2)
\psaxes[linewidth=1.25pt,Dx=5,Dy=5]{->}(0,0)(-1,-1)(4.5,2)
\psplot[plotpoints=2000,linewidth=1.25pt,linecolor=blue]{0.1}{4.5}{x ln 1.3 mul}
\uput[dr](1,0){1}\uput[d](1.8,0){$x_0$}\uput[ul](1.8,0.764){$M$}
\uput[dr](4,1.6){\blue $\mathcal{C}$}\uput[dl](0,0){O}
\psline[linestyle=dashed](1.8,0.764)(1.8,0)
\uput[d](4.4,0){$x$}\uput[l](0,1.9){$y$}
\psplotTangent{1.8}{4}{x ln 1.3 mul}
\psline(-1,0.764)(4.5,0.764)\uput[ul](0,0.764){B}\uput[ul](0, -0.5){A}
\psframe(0,0.764)(0.1,0.664)\uput[ul](3.1,1.8){$T$}
\end{pspicture*}
\end{center}

\begin{enumerate}
\setcounter{enumi}{3}
\item Démontrer que la longueur AB est égale à une constante (c'est-à-dire à un nombre qui ne dépend pas de $x_{0}$) que l'on déterminera.

\emph{Le candidat prendra soin d'expliciter sa démarche}.
\end{enumerate}

\end{document}
